\section*{Chapitre \ref{ch:rc:model-risk-and-validation} --- Risque de modèle et validation}

\begin{solution}{exo:rc:model-risk-and-validation:1}
Avec la formule voulue~: $0.10/1.05 = 9.52~\%$~; divisé par la somme~: $0.10/2.10 = 4.76~\%$, soit la moitié.
\end{solution}

\begin{solution}{exo:rc:model-risk-and-validation:2}
Score $2\times3+1+1 = 8$~: niveau~2.
\end{solution}

\begin{solution}{exo:rc:model-risk-and-validation:3}
$\P(X\ge8)$ pour une loi binomiale à 250 essais et une probabilité de 1~\%~: 0.40~\%. Improbable~: la zone jaune du
chapitre~21, avec une surcharge de fonds propres.
\end{solution}

\begin{solution}{exo:rc:model-risk-and-validation:4}
Avec une VaR divisée par deux, le taux de \omterm{def:rc:market-risk-measures:exception}{dépassements de VaR} passe de 1~\% à peut-être quelques pour cent, ce qui demande des mois de données
pour le distinguer d'une simple malchance~: à 99~\%, un modèle correct compte 2.5 dépassements par an. Les backtests détectent les
erreurs grossières lentement~; les tests d'implémentation les repèrent aussitôt.
\end{solution}

\begin{solution}{exo:rc:model-risk-and-validation:5}
Il ne teste que l'implémentation, non les hypothèses du modèle~: les deux les partagent. C'est pourtant le
test le plus exigeant des aspects numériques (pas, bornes, calibration), car tout écart est une erreur du
candidat.
\end{solution}

\begin{solution}{exo:rc:model-risk-and-validation:6}
Le modèle de VaR dans son ensemble est un modèle (il applique la théorie statistique aux données)~; le tableur qui calcule
les variations relatives fait partie de son implémentation, même si l'arithmétique élémentaire prise isolément est exclue. Son
contrôle relève de la validation du modèle et de sa gestion des changements, et non d'une seule
politique distincte pour les tableurs.
\end{solution}

\begin{solution}{exo:rc:model-risk-and-validation:7}
Avec des pas d'un vingt-quatrième d'année, 3 des 36 points à un an échouent (contre 14 avec des pas
mensuels)~; l'erreur maximale est de 1.40 points de base du notionnel, et l'erreur relative maximale, 15~\%, porte
sur une option hors de la monnaie qui vaut peu.
\end{solution}

\begin{solution}{exo:rc:model-risk-and-validation:8}
L'approbation atteste d'un examen, non de l'exactitude~; un modèle qui abaisse un chiffre pendant un dépassement de limite exige
davantage d'examen, non moins. Le modèle de 2012 a été approuvé alors qu'il était faux, et le rapport interne
de la banque a relevé que la revue n'avait pas vérifié l'implémentation.
\end{solution}

\begin{solution}{pb:rc:model-risk-and-validation:1}
\textbf{1.} $(n-o)/(n+o) = \tfrac12\,(n-o)/\bigl((n+o)/2\bigr)$.
\textbf{2.} 14.06~millions d'euros prévus, contre 7.01~millions avec l'erreur~: un rapport de 0.499.
\textbf{3.} Le portefeuille n'est pas linéaire dans les mouvements (la tranche présente une convexité), si bien que diviser les mouvements par deux ne
divise pas exactement par deux les pertes.
\textbf{4.} 12.0~: la valeur de la tranche equity varie environ douze fois autant que celle de l'indice pour une variation
de l'intensité de défaut, il faut donc douze fois le notionnel de protection sur l'indice pour le couvrir.
\textbf{5.} Les dépassements de limite disparaîtraient~: une limite fixée sur l'ancien modèle paraîtrait utilisée à moitié.
\textbf{6.} Réimplémentation indépendante de l'étape de variation relative, ou calcul à la main
d'un scénario~: aussitôt.
\textbf{7.} Plus de \omterm{def:rc:market-risk-measures:exception}{dépassements de VaR} que 1~\%, visibles seulement après des mois.
\textbf{8.} Le changement est survenu alors que les limites étaient dépassées et a fait disparaître les dépassements, ce qui a soulagé la pression
au lieu d'inciter à examiner les positions.
\textbf{9.} La documentation décrivait la formule voulue~; seule l'implémentation était fausse.
\textbf{10.} Inventaire des tableurs utilisés dans les modèles, tests indépendants de leurs formules, contrôle de version,
et édition restreinte.
\textbf{11.} Un validateur indépendant et une gestion des risques placée au-dessus du desk, avec un contrôle que les limites sont
recalibrées.
\textbf{12.} Calcul en parallèle avec l'ancien modèle, limites réduites jusqu'à ce que les résultats le confirment, et revue
de l'implémentation de chaque transformation des entrées.
\textbf{13.} Les recalibrer sur le nouveau modèle au moment du changement, afin qu'un changement de modèle ne soit pas un changement de limite.
\textbf{14.} Niveau~1~: il fixe les limites et les fonds propres d'un portefeuille vaste et complexe, aux entrées incertaines.
\textbf{15.} Compétence pour reconstruire le calcul, indépendance vis-à-vis du CIO, et autorité pour en arrêter l'usage.
\textbf{16.} Non~: ce sont les transactions et leur taille qui en sont à l'origine~; le modèle a masqué le risque et retardé la réaction.
\textbf{17.} Qu'une banque doit encore contrôler un tel calcul lorsqu'il fait partie de l'implémentation d'un modèle~;
les exclusions visent les calculs autonomes.
\textbf{18.} Tests unitaires sur des cas connus, réimplémentation indépendante des étapes clés, rapprochement des
sorties intermédiaires, et comparaison à des références sur grille.
\textbf{19.} Résultat identifié~: \emph{l'erreur de moyennage}~: diviser les variations relatives par la somme au lieu de la
moyenne les réduit de moitié et ramène la VaR à 99~\% du portefeuille de 14.06~millions d'euros à 7.01~millions, soit une réduction de 50~\%.
\textbf{20.} Le risque d'agir sur la base d'un modèle faux, mal construit ou mal utilisé.
\end{solution}

\begin{solution}{iq:rc:model-risk-and-validation:1}
Périmètre et usage~; solidité conceptuelle~; contrôles des données et de l'implémentation~; comparaison aux \omterm{def:rc:model-risk-and-validation:bench}{modèles de référence} et
analyse de sensibilité~; \omterm{def:rc:model-risk-and-validation:bench}{analyse des résultats}~; limites~; constats assortis de leur gravité, de leurs responsables
et de leurs échéances~; conditions d'usage~; conclusion du validateur.
\iqlookfor{structure et constats exploitables.}
\end{solution}

\begin{solution}{iq:rc:model-risk-and-validation:2}
Vérifier le choix du modèle au regard du produit (retour à la moyenne, smile, nombre de facteurs)~; comparer
les prix européens aux formules fermées~; comparer arbres et Monte Carlo par régression pour les options bermudéennes~; tester
les frontières d'exercice et la convergence~; vérifier la stabilité de la calibration~; comparer aux prix de marché lorsque
ils sont disponibles~; évaluer l'effet du smile avec un modèle alternatif.
\iqlookfor{\omterm{def:rc:model-risk-and-validation:bench}{modèles de référence}, convergence et modèles alternatifs.}
\end{solution}

\begin{solution}{iq:rc:model-risk-and-validation:3}
Tests unitaires à réponses connues, cas limites et symétries (parité, volatilité nulle, très dans ou très hors de la
monnaie), tests de convergence, réimplémentation indépendante, tests de régression sur les sorties stockées, et
tests de propriétés sur des grilles.
\iqlookfor{cas connus, invariants et implémentations indépendantes.}
\end{solution}

\begin{solution}{iq:rc:model-risk-and-validation:4}
Une revue critique par des personnes ayant compétence, indépendance et autorité. Elle a lieu lorsque les constats
modifient les modèles ou leur usage, lorsque les validateurs peuvent arrêter un modèle, et lorsque les désaccords sont
consignés et remontés.
\iqlookfor{les trois conditions et les éléments probants.}
\end{solution}

\begin{solution}{iq:rc:model-risk-and-validation:5}
Un nouveau modèle a abaissé la VaR pendant des dépassements de limite, a été approuvé sans vérifier son implémentation,
et comportait une erreur de tableur qui a divisé la volatilité par deux. Réimplémentation indépendante, calculs en parallèle,
recalibrage des limites lors d'un changement de modèle, et remontée des positions à l'origine des dépassements.
\iqlookfor{l'erreur, la défaillance de gouvernance et les contrôles.}
\end{solution}

\begin{solution}{iq:rc:model-risk-and-validation:6}
Évaluer le produit selon un ensemble de modèles alternatifs plausibles et de calibrations, prendre la dispersion
(ou la distance à un choix prudent) comme incertitude de modèle, et la provisionner~; réviser au fur et à mesure que
les données de consensus de marché arrivent (chapitre~27).
\iqlookfor{modèles alternatifs et provision.}
\end{solution}
